\section{Evaluating the evaluator: compositional bias and causal acceptance}
\label{sec:exec24-judge-bias}

A bank using an agent to translate business questions into causal studies also
needs a way to review the agent's work. An evaluator may compare a proposed
diagram with a narrative, check whether a report uses the supplied evidence,
or judge an explanation of uncertainty. These are useful tasks, but the
evaluator can itself ignore evidence, reward fluent unsupported answers or
react to irrelevant presentation changes. The relevant question is whether
the evaluation changes for the right reason and remains anchored to an
independently specified scientific standard.~\citep{exec24_judgebias_reference}

\citet{exec24_judgebias_original} study compositional bias in multimodal
language models used as judges. The inspected original is arXiv
2604.18164v3, dated 23 April 2026, before the research cutoff. Its title and
author order match the inherited conference-work metadata, but the final
conference PDF is not asserted to be byte-identical. The earlier v1 original
is retained separately; selected methods in the later revision, rather than
an unlabelled mixture of versions, support the discussion here.
The study concerns pointwise vision-language evaluation. Its empirical
results are not measurements of banking causal identification, diagram
validity, or the proposed library's performance.
\citep{exec24_judgebias_original}

\subsection{Two desired responses to two different changes}

The paper represents a judging input as a query $Q$, an image $I$ and a
candidate response $R$, with a score $y=f_\theta(Q,I,R)$ on a scale from one
to $S$, taking $S=10$ in its experiments. A perturbation produces
$(Q',I',R')$ and a second score $\widehat y$. The distinction between
perturbations is substantive. Removing necessary evidence or mismatching a
response with its evidence should reduce an otherwise favorable score;
changing a semantically irrelevant aspect should preserve it.
\citep{exec24_judgebias_original}

For the first type, the paper defines Bias-Deviation as
\[
 \mathrm{BD}=E\left[\frac{(y-\widehat y)_+}{y-1}\;\middle|\;y>1\right],
 \qquad (a)_+=\max(a,0).
\]
The condition excludes an original score of one because there is no room
for a further decrease. For the second type, Bias-Conformity is
\[
 \mathrm{BC}=E\left[1-
 \frac{|y-\widehat y|}{\max(y-1,S-y)}\right].
\]
The denominator normalizes by the largest available movement from the
original score. These formulas apply to different perturbation families;
using a high score from one as an answer to the other changes the target.
The study groups nine bias types into integrality, congruity and robustness.
Its method therefore requires a defensible prior judgment about whether
each perturbation changes the relevant meaning.
\citep{exec24_judgebias_original}

This requirement is especially consequential for causal diagrams. Moving a
node on the page while preserving every label and endpoint mark can be
irrelevant to the graph's meaning. Removing an arrowhead, deleting a variable
that blocks a back-door path, or moving a measurement across the decision
time is not. A generic visual transformation that is harmless in one task
can obscure an arrowhead or minus sign in another. The proposed bank
evaluation must validate the transformation against the graph or numerical
object before classifying it as meaning-preserving.
\citep{exec24_judgebias_reference}

\subsection{What the empirical design establishes}

The selected paper reports 1,804 evaluation pairs derived from 29 source
benchmarks, with four core task categories and twelve domains. It synthesizes
queries, filters them and uses human review to enforce requirements including
the need to consult the image and consistency with the specified task.
The appendix describes twelve annotators and their selection criteria. This
construction is more informative than a claim that random visual changes
automatically preserve or destroy meaning, but the resulting population is
still the curated benchmark population.~\citep{exec24_judgebias_original}

The study describes 26 judge models, including proprietary, open and
specialized critic families, and evaluates reasoning modes where available.
The main-result caption reports three experimental runs and supporting
variation across runs and samples. The stated token budget is 16,384, with
specified reasoning settings and other parameters left at model defaults.
Those operational choices are part of the inherited result. Repeated runs
help reveal variation for these cases; they do not manufacture independent
new business questions. A new comparison should preserve task identities
and report both within-case variation and variation across distinct cases.
\citep{exec24_judgebias_original,exec24_judgebias_reference}

The paper reports differing strengths across the perturbation families and
does not reduce reliability to model size or a single reasoning switch.
Its prompt interventions likewise produce different trade-offs across
models and bias types. The appropriate proposal use is to adopt the
distinction between relevant sensitivity and irrelevant instability as a
test-design principle, and to reference the observed multimodal failure
modes. Selecting a banking judge solely from its aggregate benchmark rank
would require a transfer argument and new evidence.
\citep{exec24_judgebias_original,exec24_judgebias_reference}

\subsection{Why stability is not correctness}

A constant evaluator illustrates a limitation already acknowledged in the
paper's discussion of supporting sample variance. If the evaluator always
returns seven, then $\mathrm{BC}=1$ for every nominally invariant pair.
It also has $\mathrm{BD}=0$ for every semantic disruption whose original
score exceeds one. Neither statistic by itself establishes accuracy against
a scientific oracle. In an authored four-case set containing two valid and
two invalid causal claims, accepting all four gives perfect score stability
but two false acceptances.~\citep{exec24_judgebias_original,
exec24_judgebias_reference}

Bias-Deviation has its own denominator that must remain visible. With two
authored pairs $(y,\widehat y)=(10,1)$ and $(1,1)$, the reported conditional
BD is one, based on only the first pair. The second pair is not evidence of
a successful response to the perturbation; it is outside the metric's
eligibility set. A bank comparison should publish the total number of pairs,
the number eligible for BD, their task composition and the reasons for
exclusion. It should not convert a conditional metric into an all-request
success probability.~\citep{exec24_judgebias_reference}

Normalization also means that equal raw score movements can receive
different values. A decrease from ten to eight contributes $2/9$ to BD;
a decrease from three to one contributes one. Both are two-point changes.
This is a consequence of the stated normalization, not an arithmetic
defect. It shows why base-score distributions and score-scale interpretation
matter when comparing judges. For causal acceptance, the product additionally
needs an externally defined distinction between a valid identified result,
a justified identified set, a justified refusal and an unsupported claim.
\citep{exec24_judgebias_reference}

\subsection{Abstention requires both conditional quality and coverage}

The paper's abstention-aware variant permits an N/A response, excludes that
response from the numerical metric and reports its proportion separately.
This is a clearly different estimand from a score averaged over every
submitted case. Its methodological implication is to preserve both the
quality of retained judgments and the coverage of the evaluator.
A scientific refusal is not automatically an evaluator failure, but
removing it from every denominator would obscure the service's usefulness.
\citep{exec24_judgebias_original}

Let $J$ indicate that an evaluator returns a scored judgment and $C$ indicate
that the judgment agrees with a prespecified oracle. Then
\[
 P(C\cap J)=P(C\mid J)P(J).
\]
In an authored 100-case example, 50 cases are scored and 45 of those agree
with the oracle. Conditional agreement is 90 percent, while the fraction
of all requests receiving a correct scored judgment is 45 percent. If
some unscored cases represent appropriate refusal, their contribution to
utility should be specified separately rather than labelled either correct
or incorrect by an undocumented convention.
\citep{exec24_judgebias_reference}

An economic loss makes the choice concrete. Suppose the two false
acceptances in the four-case example each expose a stipulated 250,000 USD
loss, while expert review costs 1,000 USD per case. The false-acceptance
loss is 500,000 USD and reviewing all four costs 4,000 USD. These are
hypothetical consequences, not measured bank losses. They demonstrate why
a threshold tuned to an average evaluator score need not minimize decision
loss. With a separate hypothetical expenditure of 15,000 USD for the
100-case exercise, cost is 150 USD per submission, 300 USD per scored case
and $1{,}000/3$ USD per correct scored judgment.
\citep{exec24_judgebias_reference}

\subsection{A bounded route into the proposed product}

The proposed first prototype uses public or generated causal problems with
known structural assumptions and exact refusal oracles. It constructs
paired reports and diagrams, verifies which transformations preserve the
underlying scientific object, and keeps a held-out set of business-to-causal
translations. It then evaluates sensitivity, stability, absolute correctness,
appropriate refusal, coverage and full cost together. The unit of analysis
is the logical case; for example, three attempts on each of 100 cases yield
300 attempts but only 100 distinct cases. A repeated prompt is not a new
independent validation sample.~\citep{exec24_judgebias_reference}

For integration, an evaluator request should carry the causal question,
graph artifact, evidence identities, perturbation class, rubric version,
model configuration and independent oracle where one exists. Its response
should preserve the score, explanation, abstention, failure and cost fields
through the library, CLI and MCP boundary. An LLM explanation may help a
reviewer locate a disagreement, but must not overwrite the engine's
identification or uncertainty status. Scientific owners decide the target
and acceptance loss; protocol owners preserve the evidence and denominator.
\citep{exec24_judgebias_reference}

The inspected study is a useful reference and a candidate source of test
designs. Dataset and code licensing, model access, exact artifacts and
native reproducibility still require separate inspection before direct
adoption. Porting the two explicitly defined metrics is a bounded task,
whereas claiming that a multimodal benchmark qualifies causal reasoning
requires evidence the study does not supply. Its stated limits include
vision-language, pointwise scoring and a focus on multimodal biases rather
than every social or domain-specific bias. The selected preprint supports this comparison; the final version, package execution, independent validation and applicability to banking still require evidence.
\citep{exec24_judgebias_original,exec24_judgebias_reference}
